% PART 2. CONTENT, CHAPTER 1: mô tả bài toán và yêu cầu. Nhãn PART 2 in phía trên
% chương đầu tiên và có dòng riêng trong mục lục.
\ThesisChapterNumbering
\ThesisPart{PART 2. CONTENT}
\chapter{PROBLEM DESCRIPTION}\label{chap:problem}

\section{Detailed Problem Description}
Researchers who study degrons want to know which substitution, at which position of a degron peptide,
raises or lowers the degron signal. The thesis addresses two connected problems.

\begin{itemize}
  \item \textbf{Prediction.} Given a reference peptide sequence, a position and a substituted amino acid,
        predict the effect $\Delta$PSI of the substitution. Models are trained and evaluated on
        experimental labels with protein-cluster-aware splits.
  \item \textbf{Analysis tool.} Provide a prototype that lets a user enter a peptide, view the predicted
        effect matrix, rank substitutions and export the result with provenance (model version, data
        source, supported domain).
\end{itemize}

The intended users are researchers who select mutations for follow-up experiments, and students who
explore the relation between peptide sequence and degradation signal. The workflow is: enter a valid
peptide, generate all substitutions, score them with a fixed model, display a matrix and ranking, and
export CSV/JSON. Inputs that are missing, mis-positioned or outside the supported domain must be
reported explicitly, and uncalibrated scores must not be presented as probabilities of disease.

\section{Functions of the System}
A minimal offline command-line prototype exists (Section~\ref{sec:testplan}); it covers
scoring and CSV export only, with no web interface yet. The use cases below remain the requirements for
the full prototype.

\subsection{Predict the Effect Matrix}
\begin{usecasespec}{Predict the effect matrix of a peptide}{tab:uc-predict}
  \ucname{Predict effect matrix}{UC-01}
  \ucactor{User}{Important}{Simple}
  \ucfull{Stakeholders and interests:}{Researcher: wants predicted effect of every substitution of a peptide.}
  \ucfull{Summary description:}{The user enters a peptide sequence; the system returns an $L\times 20$ matrix of predicted $\Delta$PSI values.}
  \ucfull{Relationships:}{Association: User \newline Include: Validate input \newline Extend: Compare with measured data \newline Generalization: ---}
  \ucfull{Event handling:}{1.~The user enters a peptide. \newline 2.~The system validates length, alphabet and domain. \newline 3.~The system scores all substitutions with the fixed model. \newline 4.~The system displays the matrix.}
  \ucfull{Alternative / exceptional flows:}{Invalid or out-of-domain input: the system shows an explicit message and does not return scores.}
  \ucfull{Expected outcome:}{A matrix of predicted effects with model version and supported domain.}
\end{usecasespec}

\subsection{Rank Substitutions}
\begin{usecasespec}{Rank substitutions}{tab:uc-rank}
  \ucname{Rank substitutions}{UC-02}
  \ucactor{User}{Important}{Simple}
  \ucfull{Stakeholders and interests:}{Researcher: wants a short list of substitutions to test.}
  \ucfull{Summary description:}{The system orders substitutions by predicted effect.}
  \ucfull{Relationships:}{Association: User \newline Include: Predict effect matrix \newline Extend: --- \newline Generalization: ---}
  \ucfull{Event handling:}{1.~The user opens the ranking of a predicted peptide. \newline 2.~The system sorts substitutions by predicted effect. \newline 3.~The user filters by position or amino acid.}
  \ucfull{Alternative / exceptional flows:}{No prediction available: the system asks the user to run UC-01 first.}
  \ucfull{Expected outcome:}{A ranked table of substitutions.}
\end{usecasespec}

\subsection{Export Results with Provenance}
\begin{usecasespec}{Export results with provenance}{tab:uc-export}
  \ucname{Export results}{UC-03}
  \ucactor{User}{Medium}{Simple}
  \ucfull{Stakeholders and interests:}{Researcher: wants results that can be reused and traced.}
  \ucfull{Summary description:}{The system exports matrix and ranking with model version, data source and supported domain.}
  \ucfull{Relationships:}{Association: User \newline Include: Predict effect matrix \newline Extend: --- \newline Generalization: ---}
  \ucfull{Event handling:}{1.~The user selects an export format (CSV or JSON). \newline 2.~The system writes the file with provenance fields.}
  \ucfull{Alternative / exceptional flows:}{Export fails: the system reports the error and keeps the result on screen.}
  \ucfull{Expected outcome:}{A file containing scores and provenance.}
\end{usecasespec}

\section{Non-Functional Requirements}
These are targets for the prototype; \textbf{[PENDING: concrete thresholds to be fixed with the system design.]}
\subsection{Performance Requirements}
Scoring a single peptide should finish in interactive time on a standard laptop (target to be measured).
\subsection{Information Safety Requirements}
Only public data are used; terms of use of each source are recorded before redistribution.
\subsection{Security Requirements}
Input is validated; user-supplied sequences are not executed or stored beyond the session unless stated.
\subsection{Reliability Requirements}
The same input and model version must give the same scores; scores shown in the interface must match the
research runner within numerical tolerance.
\subsection{Maintainability Requirements}
Models and data are versioned; experiments can be re-run from documented commands.

\section{Implementation and Design Constraints}
\subsection{Implementation Constraints}
The work is done by one student within the thesis course schedule; no new wet-lab experiment is performed;
computation runs on CPU-class hardware.
\subsection{Design Constraints}
Scores are reported for the assay and domain on which the model was evaluated; no claim is made on
disease causation or treatment.
